\section{Introduction}\label{introduction}

Quantum change-point inference asks where an ordered source changes from one
state-generating process to another. Unlike a classical observer, the receiver
may retain the systems and measure the entire sequence collectively.
Early Bayesian work considered unknown pure states and two possible change
locations \cite{akimoto2011}. For a single permanent change, collective
minimum-error discrimination has a nonzero asymptotic success probability and
can outperform online local strategies \cite{sentis2016}; exact and online
unambiguous-identification protocols were subsequently developed in
Refs.~\cite{sentis2017,sentis2018}. An adaptive Bayesian online protocol was
also demonstrated experimentally with photonic state sequences
\cite{yu2018}. Separate formulations optimize different operational risks.
Quickest and sequential quantum change detection trade false-alarm control
against detection delay \cite{fanizza2023,zecchin2026}. Hamiltonian change
points are finite-candidate quantum process-discrimination problems under a
minimum-error criterion \cite{nakahira2023}, whereas the 2026 channel result
studies unambiguous process discrimination and hence permits an inconclusive
outcome \cite{nakahira2026}. The literature also includes general Gram-matrix
semidefinite programs, multiple changes, isolated anomalies, and boundaries
between quantum domains
\cite{mohan2023,skotiniotis2024,llorens2024,llorens2025}.

Here the change is transient: the source returns to its reference state after
one nonempty contiguous anomalous block. We refer to this
interval-constrained family as a \emph{quantum change interval}. For a
sequence of length \(n\) and an interval \(I=[a,b]\), the corresponding
hypothesis is

\[
|\Phi_{a,b}\rangle
=|0\rangle^{\otimes(a-1)}|\psi\rangle^{\otimes(b-a+1)}
|0\rangle^{\otimes(n-b)}.
\]

When the length \(i=b-a+1\) is known, only the translation \(a\) is inferred;
when the length is unknown, both endpoints are inferred.
Figure~\ref{fig:model-geometry} summarizes this returning geometry and the
collective exact-localization task.

The model is related to established quantum tasks through several precise
restrictions. The
boundary slice \(b=n\) is the permanent quantum change-point family
\cite{sentis2016,sentis2017,sentis2018}. A general multiple-change-point
framework contains it as the constrained pattern
\(\lvert0\rangle\to\lvert\psi\rangle\to\lvert0\rangle\)
\cite{mohan2023}. For fixed length \(i\), the candidates form the strict
subensemble of the \(k=i\) multi-anomaly problem in which the anomalous sites
are consecutive; the unknown-length family is the union of these contiguous
subensembles over \(i\) \cite{llorens2024}. Quantum edge detection instead
locates a single boundary between domains whose pure states are unknown and
proves an SRM-based protocol asymptotically optimal
\cite{llorens2025}. It is therefore a single-boundary geometric analogue, not
a prior-preserving supermodel of the present task. Here two endpoints are tied
to one interval, the two local states are calibrated and known, and both
exterior domains are fixed to the same reference state \(\lvert0\rangle\).

The unknown-length problem has a specific obstruction that is hidden by the
two-endpoint notation. For intersecting intervals, the overlap kernel depends
only on the endpoint displacement and agrees with a separable two-dimensional
Toeplitz kernel. For separated intervals, however, the physical overlap
contains an additional correction determined by the gap between the two
intervals. Thus the exact physical Gram matrix is neither a global
two-dimensional Toeplitz compression nor an element of a transitive
association scheme. A direct multilevel Szeg{\H o} theorem therefore does not
solve the problem. The proof deletes a short-interval sector whose prior mass
vanishes, compares the retained long-interval Gram matrix in operator norm
with a separable convolution kernel on a triangular Følner sequence, and
obtains matching SRM and optimum limits through local-diagonal functional
calculus and a Gram upper bound. A prior-mass sandwich transfers the optimum
to all intervals, while principal-submatrix interlacing, including the
essential retained-family factor, transfers the square-root trace and hence
the SRM of the full physical family.

\begin{figure*}[t]
  \centering
  \includegraphics[width=\textwidth]{figures/figure1_model_geometry.pdf}
  \caption{\textbf{Returning quantum change interval and collective localization.}
  A source emits the reference state \(\lvert0\rangle\) except on a nonempty
  contiguous interval \(I=[a,b]\), where it emits the known anomalous state
  \(\lvert\psi\rangle\). The entry and return boundaries are \(a\) and
  \(b+1\), respectively. For known length, \(i=b-a+1\) is fixed and only
  \(a\) is inferred; for unknown length, both endpoints are inferred. An
  arbitrary collective POVM may act on all \(n\) systems and returns an exact
  interval estimate \(\widehat I\).}
  \label{fig:model-geometry}
\end{figure*}

Throughout this paper, \emph{exact localization} refers to the zero--one
reward: an outcome is successful only when the reported interval
\(\widehat I\) equals the true interval \(I\). The minimum-error POVM has one
reported interval for every outcome and may make an incorrect report. This is
different from \emph{unambiguous exact identification}, where every
conclusive report must be correct, the cross-errors are constrained to zero,
and an additional inconclusive outcome is allowed
\cite{sentis2017,llorens2024}. The present paper studies the former criterion,
including when a no-anomaly hypothesis is added.

One operational setting is a batch of quantum probes or carriers emitted by
a device during a transient coherent disturbance. The model assumes
independent single-system outputs, a calibrated pure anomalous state
\(\lvert\psi\rangle\), an abrupt switch from \(\lvert0\rangle\), a state that
remains constant throughout the anomalous block, and a perfect return to
\(\lvert0\rangle\). A receiver may retain the complete batch coherently and
apply an unrestricted global POVM to estimate the onset and duration. We treat
this as an ultimate collective-measurement benchmark. Correlated or mixed
outputs, gradual switching, imperfect recovery, unknown fault states, and
restricted storage or readout architectures require different models.

The same returning geometry appears in classical statistics. Epidemic
alternatives and changed-segment models place a temporary distributional
change on an unknown contiguous
interval \cite{levin1985,yao1993}, while scan methods cover spatial or
graph-structured changed regions \cite{walther2010,chen2015}. Those works
optimize classical test statistics and type-I-error or power criteria. Here
the observations are noncommuting quantum hypotheses, arbitrary collective
POVMs are allowed, and the target is the exact Bayes localization probability.
Table~\ref{tab:literature} summarizes the resulting distinctions.

\begin{table*}[t]
\centering
\caption{Theorem-level comparison and structural relations between quantum
change intervals and neighboring quantum and classical models.}
\label{tab:literature}
\scriptsize
\begin{tabularx}{\textwidth}{@{}
>{\raggedright\arraybackslash}p{0.15\textwidth}
>{\raggedright\arraybackslash}p{0.21\textwidth}
>{\raggedright\arraybackslash}p{0.30\textwidth}
>{\raggedright\arraybackslash}X@{}}
\toprule
Work & Relation to change intervals & Established result & Difference addressed here \\
\midrule
Sentís et al.~\cite{sentis2016,sentis2017,sentis2018}
& Exactly the boundary slice \(b=n\)
& Minimum-error permanent-change limit; separate finite-size
unambiguous-identification laws and online protocols
& A returning interval couples entry and return boundaries under the
minimum-error criterion \\
Mohan--Sikora--Upadhyay \cite{mohan2023}
& Contains this model as a constrained two-change-point instance
& Their Theorem~2.1 gives an exact Gram-only SDP for general Bayes rewards;
Appendix~D gives a sectorwise two-change Gram kernel and finite SDP/heuristic
calculations
& No closed \(n\to\infty\) interval law or SRM-optimality theorem for the
returning subfamily \\
Quantum multi-anomaly detection \cite{llorens2024}
& For fixed \(i\), contains the interval candidates as the consecutive \(k=i\) subensemble
& Johnson symmetry gives exact finite-\(n\) SRM optimality and a closed
spectral formula; for fixed \(k\),
\(P_{\rm opt}=(1-c^2)^k+O(n^{-1/2})\), plus a zero-error program
& The consecutive-subset restriction destroys Johnson transitivity and
requires Toeplitz/Følner reductions \\
Quantum edge detection \cite{llorens2025}
& Single-boundary geometric analogue, not a prior-preserving containment relation
& Schur--Weyl and SDP analysis for unknown pure-state domains; an SRM protocol
is asymptotically optimal
& Here two tied endpoints delimit one interval, the states are known, the
exterior domains equal \(\lvert0\rangle\), and the full physical Gram kernel
requires a non-Toeplitz-to-Følner transfer \\
Epidemic/changed-interval statistics \cite{levin1985,yao1993,walther2010,chen2015}
& Direct classical analogue of the returning interval
& Scan/CUSUM detection, significance, and power guarantees for commuting data
& The quantum problem optimizes exact Bayes localization over collective POVMs \\
\bottomrule
\end{tabularx}
\end{table*}

We use \emph{quantum change interval} as a descriptive name for this returning
contiguous ensemble. It remains a restricted instance of multiple quantum
change points and anomalous-site identification, and is distinct from the
unknown-state boundary problem in quantum edge detection. Standard Toeplitz
and circulant asymptotics, including multilevel spectral-distribution
theorems, provide the matrix-analytic background
\cite{gray2006,tilli1998,tyrtyshnikov1998}. Asymptotic SRM optimality by
itself is not unique to this setting: it is already established for quantum
edge detection \cite{llorens2025}. The contribution here is the combination
of exact asymptotic Bayes laws for the interval ensemble, a Toeplitz/Følner
reduction that handles a physical kernel which is not globally
two-dimensional Toeplitz, and a transfer proving asymptotic optimality for the
SRM of the full, untruncated physical family. More precisely:

\begin{enumerate}
\def\labelenumi{\arabic{enumi}.}
\tightlist
\item
  For fixed known length \(i\), the optimal and SRM success probabilities converge to a one-dimensional Toeplitz integral, with an \(O(N^{-1/2})\) bound on their difference.
\item
  If \(i=i_n\to\infty\) and the number \(N_n=n-i_n+1\) of admissible translations also diverges, both probabilities converge to \(p_1(c^2)\), with no assumption on \(i_n/N_n\).
\item
  Under the prior uniform on all nonempty intervals, both probabilities
  converge to \(p_1(c)^2\). This law is tied to that interval prior and is not
  asserted for arbitrary priors on the unknown length. The change from
  \(p_1(c^2)\) to \(p_1(c)^2\) reflects both the replacement of one translation
  coordinate by two endpoint coordinates and the change of the elementary
  overlap parameter from \(c^2\) to \(c\).
\item
  If the no-anomaly prior \(\pi_0\in(0,1)\) is fixed and the conditional localization limit is \(L\), the joint detection-and-localization success probability converges to \(\pi_0+(1-\pi_0)L\).
\end{enumerate}

The endpoint cases \(c=0,1\) are treated separately. Throughout, the source
outputs are independent, the switching is abrupt, and the anomalous block has
a fixed calibrated pure state before a perfect return to \(\lvert0\rangle\).
The overlap \(c\) is independent of \(n\); when \(H_0\) is present,
\(\pi_0\in(0,1)\) is likewise independent of \(n\). The location or interval
has the stated uniform prior, and coherent storage followed by an arbitrary
collective POVM is allowed. We consider one anomalous interval; local
asymptotic alternatives \(c=c_n\to1\), unknown states, measurement
restrictions, and multiple disjoint intervals are outside the scope of the
results below.

\section{Model and discrimination preliminaries}\label{model-and-discrimination-preliminaries}

\subsection{Phase convention and Gram matrices}\label{phase-convention-and-gram-matrices}

Write

\[
\langle0|\psi\rangle=ce^{i\varphi},\qquad 0\le c\le1,
\qquad r=c^2.
\]

Replacing \(|\psi\rangle\) by \(e^{-i\varphi}|\psi\rangle\) changes each candidate vector only by a hypothesis-dependent global phase. Equivalently, the original and phase-fixed Gram matrices are related by a diagonal unitary conjugation. Minimum-error and SRM success probabilities are unchanged. We therefore use

\[
\langle0|\psi\rangle=c\ge0
\]

without loss of generality.

Unless a theorem explicitly introduces a sequence \(i=i_n\), every
asymptotic statement is taken as \(n\to\infty\) with \(c\) (and hence
\(r=c^2\)) fixed. In Section~\ref{joint-detection-and-exact-localization},
the no-anomaly prior \(\pi_0\in(0,1)\) is also fixed rather than scaled with
\(n\).

For an ensemble of \(M\) normalized pure states \(\{|\phi_j\rangle\}_{j=1}^M\) with equal priors, let

\[
G_{jk}=\langle\phi_j|\phi_k\rangle.
\]

We write \(P_{\rm opt}(G)\) for the minimum-error optimum over all POVMs
\cite{helstrom1976,holevo1973,yuen1975,barnett2009} and

\[
P_{\rm SRM}(G)=\frac1M\sum_{j=1}^M(\sqrt G)_{jj}^2
\]

for the square-root-measurement success probability
\cite{hausladen1994,eldar2001,dallapozza2015}. If \(G\) is singular,
\(G^{-1/2}\) is understood on the support of \(G\), and the POVM is completed
arbitrarily on its orthogonal complement; this does not change any ensemble
success probability. The elementary trace bound

\begin{equation}\label{eq:trace-srm-bound}
\left(\frac{\operatorname{tr}\sqrt G}{M}\right)^2
\le P_{\rm SRM}(G)\le P_{\rm opt}(G)
\end{equation}

will be used repeatedly.

\subsection{Stability under Gram perturbations}\label{stability-under-gram-perturbations}

\begin{lemma}[Operator-norm Gram stability]\label{lem:gram-stability}
Let \(G,H\in\mathbb C^{M\times M}\) be positive semidefinite with unit diagonal. Then

\begin{equation}\label{eq:opt-gram-stability}
|P_{\rm opt}(G)-P_{\rm opt}(H)|
\le \|\sqrt G-\sqrt H\|_{\rm op}
\le \|G-H\|_{\rm op}^{1/2}
\end{equation}

and

\begin{align}\label{eq:srm-gram-stability}
&|P_{\rm SRM}(G)-P_{\rm SRM}(H)| \notag\\
&\qquad\le2\|\sqrt G-\sqrt H\|_{\rm op}
\le2\|G-H\|_{\rm op}^{1/2}.
\end{align}

If \(G,H\ge \mu I\) for some \(\mu>0\), then

\begin{equation}\label{eq:sqrt-lipschitz}
\|\sqrt G-\sqrt H\|_{\rm op}
\le \frac{\|G-H\|_{\rm op}}{2\sqrt\mu}.
\end{equation}
\end{lemma}

\begin{proof}
Use the canonical Gram realizations \(|x_j\rangle=\sqrt G e_j\) and
\(|y_j\rangle=\sqrt H e_j\). Their norms are one because the Gram matrices
have unit diagonal, and the trace distance between the corresponding pure
states is at most \(\|x_j-y_j\|\). Apply an optimal POVM for either ensemble
to the other and average over \(j\) to obtain the first inequality in
\eqref{eq:opt-gram-stability}. The second is the operator Hölder inequality
for positive square roots. Since
\(0\le(\sqrt G)_{jj},(\sqrt H)_{jj}\le1\), the scalar bound
\(|a^2-b^2|\le2|a-b|\) gives \eqref{eq:srm-gram-stability}. For the final
claim, \(X=\sqrt G-\sqrt H\) solves

\[
\sqrt G X+X\sqrt H=G-H.
\]

The inverse Sylvester operator has norm at most \((2\sqrt\mu)^{-1}\), which proves \eqref{eq:sqrt-lipschitz}.
\end{proof}

\section{A known interval length}\label{a-known-interval-length}

Let the anomalous interval have known length \(i\ge1\). There are

\[
N=n-i+1
\]

admissible starting points, and the equal-prior candidates are

\begin{align*}
|\Phi_a^{(i)}\rangle
={}&|0\rangle^{\otimes(a-1)}|\psi\rangle^{\otimes i}
|0\rangle^{\otimes(n-a-i+1)}, \\
&a=1,\ldots,N.
\end{align*}

Their exact Gram matrix is

\begin{equation}\label{eq:fixed-gram}
(G_{N,i})_{ab}=r^{\min(|a-b|,i)}.
\end{equation}

Define

\begin{equation}\label{eq:fixed-symbol}
f_{i,r}(\theta)
=1-r^i+2\sum_{d=1}^{i-1}(r^d-r^i)\cos(d\theta).
\end{equation}

\subsection{Fixed-length limit}\label{fixed-length-limit}

\begin{theorem}[Fixed-length localization]\label{thm:fixed-length}
For every fixed integer \(i\ge1\) and \(0\le r\le1\),

\begin{equation}\label{eq:fixed-limit}
\lim_{N\to\infty}P_{\rm opt}(G_{N,i})
=\lim_{N\to\infty}P_{\rm SRM}(G_{N,i})
=P_\infty(i,r),
\end{equation}

where

\begin{equation}\label{eq:fixed-integral}
P_\infty(i,r)
=\left[
\frac1{2\pi}\int_0^{2\pi}\sqrt{f_{i,r}(\theta)}\,d\theta
\right]^2.
\end{equation}

For fixed \(i\) and \(0<r<1\),

\begin{equation}\label{eq:fixed-gap-rate}
0\le P_{\rm opt}(G_{N,i})-P_{\rm SRM}(G_{N,i})
=O_{i,r}(N^{-1/2}).
\end{equation}
\end{theorem}

\begin{proof}
Put \(q=r^i\) and decompose

\[
G_{N,i}=qJ_N+B_{N,i},
\]

where \(B_{N,i}=T_N(f_{i,r})\) has entries

\[
(B_{N,i})_{ab}=
\begin{cases}
r^{|a-b|}-r^i,&|a-b|<i,\\
0,&|a-b|\ge i.
\end{cases}
\]

Choose a local basis with \(|\psi\rangle=\sqrt r\,|0\rangle+\sqrt{1-r}\,|1\rangle\). After removing the common all-zero component, each candidate has a basis component, with squared amplitude \((1-r)^i\), that occurs in no other candidate. Hence

\begin{equation}\label{eq:residual-spectral-gap}
B_{N,i}\ge(1-r)^iI_N.
\end{equation}

The corresponding infinite residual convolution operator has symbol \(f_{i,r}\), so \eqref{eq:residual-spectral-gap}, applied to arbitrarily long finitely supported vectors, implies

\begin{equation}\label{eq:symbol-spectral-gap}
f_{i,r}(\theta)\ge(1-r)^i.
\end{equation}

For \(N\ge2i\), let \(\widetilde B_{N,i}\) be the circulant matrix with the same local coefficients and set

\[
C_{N,i}=r^iJ_N+\widetilde B_{N,i}.
\]

Its eigenvalues away from the uniform Fourier mode are the sampled values \(f_{i,r}(2\pi k/N)\); the uniform mode receives the additional positive rank-one contribution. Thus \(C_{N,i}\) is a positive definite, unit-diagonal circulant Gram matrix. The matrices \(B_{N,i}\) and \(\widetilde B_{N,i}\) differ only in the upper-right and lower-left wrap-around corner blocks.

To make the boundary estimate explicit, set

\begin{align*}
b_d&=r^d-r^i,\\
E_{N,d}&=\sum_{a=1}^{d}e_a e_{N-d+a}^*,
\qquad 1\le d<i.
\end{align*}

Here \(E_{N,d}\) is a rank-\(d\) partial permutation from the last \(d\)
coordinates to the first \(d\) coordinates, and

\begin{equation}\label{eq:corner-partial-permutation}
\widetilde B_{N,i}-B_{N,i}
=\sum_{d=1}^{i-1}b_d(E_{N,d}+E_{N,d}^*).
\end{equation}

For \(N\ge2i\), the initial and terminal coordinate subspaces in
\eqref{eq:corner-partial-permutation} are disjoint. Consequently
\(\|E_{N,d}+E_{N,d}^*\|_1=2d\). Moreover, the whole difference maps between
the span of the first \(i-1\) coordinates and the span of the last \(i-1\)
coordinates. The triangle inequality and this two-corner support therefore give

\begin{equation}\label{eq:circulant-trace-perturbation}
\|G_{N,i}-C_{N,i}\|_1
\le K_{i,r}:=2\sum_{d=1}^{i-1}d(r^d-r^i),
\end{equation}

\begin{equation}\label{eq:circulant-rank-perturbation}
\operatorname{rank}(G_{N,i}-C_{N,i})\le2(i-1).
\end{equation}

The Sylvester integral gives

\begin{equation}\label{eq:sqrt-trace-perturbation}
\|\sqrt{G_{N,i}}-\sqrt{C_{N,i}}\|_1
\le\frac{K_{i,r}}{2(1-r)^{i/2}}=O_{i,r}(1).
\end{equation}

Because \(C_{N,i}\) and \(\sqrt{C_{N,i}}\) are circulant, the diagonal of \(\sqrt{C_{N,i}}\) is constant. If

\[
\beta_N=\frac1N\operatorname{tr}\sqrt{C_{N,i}},
\]

then the SRM value is \(\beta_N^2\). The Yuen--Kennedy--Lax dual operator

\[
Y=\frac{\beta_N}{N}\sqrt{C_{N,i}}
\]

is feasible because, with \(S=\sqrt{C_{N,i}}\),

\[
\beta_NS-Se_je_j^*S\ge0.
\]

Indeed, since \(S\succ0\), congruence by \(S^{-1/2}\) reduces this to

\[
\beta_N I-S^{1/2}e_je_j^*S^{1/2}\ge0.
\]

The rank-one term has unique nonzero eigenvalue \(e_j^*Se_j=\beta_N\), because the diagonal of \(S\) is constant. Since \(\operatorname{tr}Y=\beta_N^2\), the SRM is exactly optimal for the circulant model.

Let

\[
\tau_N=\|\sqrt{G_{N,i}}-\sqrt{C_{N,i}}\|_1.
\]

Put \(\Delta_N=\sqrt{G_{N,i}}-\sqrt{C_{N,i}}\). The canonical Gram realizations, followed by Cauchy--Schwarz and the monotonicity of Schatten norms, give the explicit chain

\begin{align}\label{eq:opt-circulant-transfer}
|P_{\rm opt}(G_{N,i})-P_{\rm opt}(C_{N,i})|
&\le\frac1N\sum_{j=1}^N\|\Delta_Ne_j\|_2 \notag\\
&\le\frac1{\sqrt N}
\left(\sum_{j=1}^N\|\Delta_Ne_j\|_2^2\right)^{1/2} \notag\\
&=\frac{\|\Delta_N\|_{\rm HS}}{\sqrt N}
\notag\\
&\le\frac{\|\Delta_N\|_1}{\sqrt N}
=\frac{\tau_N}{\sqrt N}.
\end{align}

while the SRM diagonal formula gives

\[
|P_{\rm SRM}(G_{N,i})-P_{\rm SRM}(C_{N,i})|
\le\frac{2\tau_N}{N}.
\]

Together with \eqref{eq:sqrt-trace-perturbation}, this proves \eqref{eq:fixed-gap-rate}. Finally, the sampled eigenvalues form a Riemann sum. The rank-one uniform-mode correction changes \(N^{-1}\operatorname{tr}\sqrt C\) by at most \(O(N^{-1/2})\), and therefore

\[
\frac1N\operatorname{tr}\sqrt{C_{N,i}}
\longrightarrow
\frac1{2\pi}\int_0^{2\pi}\sqrt{f_{i,r}(\theta)}\,d\theta.
\]

This proves \eqref{eq:fixed-limit}--\eqref{eq:fixed-integral} for \(0<r<1\). The cases \(r=0,1\) are handled below.
\end{proof}

\begin{corollary}[Residual Toeplitz family]\label{cor:residual-toeplitz}
Fix \(i\ge1\) and \(0\le r<1\), put \(q=r^i\), and let

\[
R_{N,i}:=\frac{G_{N,i}-qJ_N}{1-q}
=T_N\!\left(\frac{f_{i,r}}{1-r^i}\right).
\]

Then \(R_{N,i}\) is the Gram matrix of normalized residual states and

\begin{equation}\label{eq:residual-toeplitz-limit}
\lim_{N\to\infty}P_{\rm opt}(R_{N,i})
=\lim_{N\to\infty}P_{\rm SRM}(R_{N,i})
=\frac{P_\infty(i,r)}{1-r^i}.
\end{equation}

Equivalently,
\((1-r^i)P_{\rm opt}(R_{N,i})\to P_\infty(i,r)\), and the same statement
holds with \(P_{\rm SRM}\).
\end{corollary}

\begin{proof}
The diagonal of \(G_{N,i}-qJ_N\) is \(1-q\), and
\eqref{eq:residual-spectral-gap} shows that it is positive definite. Thus
\(R_{N,i}\) is a unit-diagonal Gram matrix. Its Toeplitz symbol is
\(h_{i,r}=f_{i,r}/(1-r^i)\), which is bounded above and below away from zero.
The one-dimensional instance of Lemma~\ref{lem:folner-diagonal}, followed by
Lemma~\ref{lem:gram-upper-bound} and \eqref{eq:trace-srm-bound}, gives the common
limit

\[
\left[\frac1{2\pi}\int_0^{2\pi}\sqrt{h_{i,r}(\theta)}\,d\theta\right]^2
=\frac{P_\infty(i,r)}{1-r^i}.
\]
\end{proof}

\subsection{Exact and special-function cases}\label{exact-and-special-function-cases}

\begin{corollary}[One-site interval, finite size]\label{cor:one-site}
For \(i=1\),

\[
G_{N,1}=(1-r)I_N+rJ_N
\]

is permutation invariant. Its positive square root has constant diagonal, so the preceding dual certificate proves finite-size SRM optimality. Since its eigenvalues are \(1+(N-1)r\) once and \(1-r\) with multiplicity \(N-1\),

\begin{align}\label{eq:one-site-finite}
P_{\rm opt}(G_{N,1})
&=P_{\rm SRM}(G_{N,1}) \notag\\
&=\frac{\bigl[\sqrt{1+(N-1)r}
+(N-1)\sqrt{1-r}\bigr]^2}{N^2}.
\end{align}

In particular, \(P_\infty(1,r)=1-r\).

This is the \(k=1\) specialization of the known single-anomaly result in
Refs.~\cite{skotiniotis2024,llorens2024}, after identifying their anomaly
overlap parameter with \(r=c^2\). We include it to show that the interval
family recovers that established finite-size formula at \(i=1\), not as a
new single-anomaly theorem.
\end{corollary}

\begin{corollary}[Two-site interval]\label{cor:two-site}
With

\[
E(m)=\int_0^{\pi/2}\sqrt{1-m\sin^2u}\,du,
\]

one has

\begin{equation}\label{eq:two-site-elliptic}
P_\infty(2,r)
=\frac4{\pi^2}(1-r)(1+3r)
E^2\!\left(\frac{4r}{1+3r}\right).
\end{equation}
\end{corollary}

\begin{proof}[Derivation]
Equation~\eqref{eq:fixed-symbol} becomes

\begin{align*}
f_{2,r}(\theta)
&=(1-r)(1+r+2r\cos\theta) \\
&=(1-r)(1+3r-4r\sin^2(\theta/2)).
\end{align*}

With \(u=\theta/2\) and symmetry about \(\pi/2\),

\begin{align*}
&\frac1{2\pi}\int_0^{2\pi}\sqrt{f_{2,r}(\theta)}\,d\theta \\
&\qquad=\frac2\pi\sqrt{(1-r)(1+3r)}
E\!\left(\frac{4r}{1+3r}\right).
\end{align*}

Squaring proves \eqref{eq:two-site-elliptic}.
\end{proof}

\subsection{Endpoint parameters}\label{endpoint-parameters}

If \(r=0\), then \(G_{N,i}=I_N\) and both success probabilities equal one. If \(r=1\), then \(G_{N,i}=J_N\), all hypotheses represent the same ray, and both probabilities equal \(1/N\). These values agree with the continuous limiting formulas. All spectral-gap arguments above are used only for \(0<r<1\).

\section{A growing known length}\label{a-growing-known-length}

Define the Poisson kernel

\[
g_r(\theta)=\sum_{d\in\mathbb Z}r^{|d|}e^{id\theta}
=\frac{1-r^2}{1-2r\cos\theta+r^2}
\]

and

\begin{equation}\label{eq:p1-definition}
p_1(r)=\left[\frac1{2\pi}\int_0^{2\pi}\sqrt{g_r(\theta)}\,d\theta\right]^2
=\frac{4(1-r^2)}{\pi^2}K^2(r^2),
\end{equation}

where \(K\) is the complete elliptic integral of the first kind in the parameter convention and \(p_1(1)=0\) by continuity. The last equality follows from the standard Landen transformation; the integral representation is the definition used in all proofs.

\begin{proposition}[One-change Toeplitz benchmark]\label{prop:one-change-benchmark}
For

\[
Q_N=(r^{|a-b|})_{a,b=1}^N,
\qquad 0\le r<1,
\]

one has

\begin{align}\label{eq:one-change-limit}
P_{\rm opt}(Q_N),P_{\rm SRM}(Q_N)
&\longrightarrow p_1(r), \notag\\
\frac1N\operatorname{tr}\sqrt{Q_N}
&\longrightarrow\sqrt{p_1(r)}.
\end{align}
\end{proposition}

\begin{proof}
The matrix \(Q_N\) is the compression of the one-dimensional convolution operator with kernel \(r^{|d|}\), whose symbol is \(g_r\) and is bounded above and below away from zero. Lemma~\ref{lem:folner-diagonal}, with \(d=1\) and \(\phi(t)=\sqrt t\), gives convergence in mean of the square-root diagonal to \((2\pi)^{-1}\int\sqrt{g_r}=\sqrt{p_1(r)}\). Lemma~\ref{lem:gram-upper-bound} supplies the matching minimum-error upper bound, while \eqref{eq:trace-srm-bound} supplies the lower bound and the SRM squeeze. This also recovers the benchmark of Ref. \cite{sentis2016}.
\end{proof}

\begin{theorem}[Long known interval]\label{thm:long-known}
Let \(i=i_n\) and \(N_n=n-i_n+1\). If

\begin{equation}\label{eq:long-known-conditions}
i_n\to\infty,\qquad N_n\to\infty,
\end{equation}

then, for every fixed \(0\le r\le1\),

\begin{equation}\label{eq:long-known-limit}
P_{\rm opt}(G_{N_n,i_n}),
P_{\rm SRM}(G_{N_n,i_n})
\longrightarrow p_1(r)=p_1(c^2).
\end{equation}

No condition on \(i_n/N_n\) is required.
The convergence is pointwise in fixed \(r\). No finite-size rate or
uniformity as \(r\uparrow1\) is claimed.
\end{theorem}

\begin{proof}
The endpoint values \(r=0,1\) follow from the matrices described in
Section~\ref{endpoint-parameters}. We therefore assume \(0<r<1\) and put
\(q_i=r^i\). Separating the common all-zero component gives

\[
|\Phi_a^{(i)}\rangle
=\sqrt{q_i}|\Omega\rangle
+\sqrt{1-q_i}|\xi_a^{(i)}\rangle,
\qquad \langle\Omega|\xi_a^{(i)}\rangle=0.
\]

The normalized residual Gram matrix is

\begin{equation}\label{eq:residual-gram}
(R_{N,i})_{ab}=
\begin{cases}
\dfrac{r^{|a-b|}-r^i}{1-r^i},&|a-b|<i,\\
0,&|a-b|\ge i.
\end{cases}
\end{equation}

Its symbol

\[
h_{i,r}(\theta)
=1+2\sum_{d=1}^{i-1}\frac{r^d-r^i}{1-r^i}\cos(d\theta)
\]

satisfies

\begin{equation}\label{eq:residual-symbol-convergence}
\|h_{i,r}-g_r\|_\infty
\le\frac{2(i-1)r^i}{1-r^i}+\frac{2r^i}{1-r}
=:\varepsilon_i(r)\longrightarrow0.
\end{equation}

Let \(Q_N=(r^{|a-b|})_{a,b=1}^N\). Toeplitz compression gives

\begin{equation}\label{eq:residual-operator-convergence}
\|R_{N,i}-Q_N\|_{\rm op}\le\varepsilon_i(r).
\end{equation}

Since \(g_r\ge(1-r)/(1+r)>0\), both matrices have a common positive
spectral gap for large \(i\). Lemma~\ref{lem:gram-stability} therefore
transfers the minimum-error optimum uniformly in \(N\). The trace distance
between each original state and its residual state is also
\(\sqrt{q_i}=r^{i/2}\), so

\[
|P_{\rm opt}(G_{N,i})-P_{\rm opt}(R_{N,i})|
\le r^{i/2}.
\]

The single-change-point Toeplitz limit for \(Q_N\) therefore proves

\[
P_{\rm opt}(G_{N,i})\to p_1(r).
\]

For the SRM, the rank-one common component requires a separate trace estimate.
Since

\[
G_{N,i}=q_iJ_N+(1-q_i)R_{N,i},
\]

operator monotonicity and the Rotfel'd trace inequality give the exact sandwich

\begin{align}\label{eq:rank-one-sandwich}
\sqrt{1-q_i}\,\operatorname{tr}\sqrt{R_{N,i}}
&\le\operatorname{tr}\sqrt{G_{N,i}} \notag\\
&\le\sqrt{1-q_i}\,\operatorname{tr}\sqrt{R_{N,i}}
+\sqrt{q_iN}.
\end{align}

Equation~\eqref{eq:residual-operator-convergence}, the common spectral gap, and the Toeplitz limit imply

\[
\frac1N\operatorname{tr}\sqrt{R_{N,i}}
\longrightarrow\sqrt{p_1(r)}.
\]

After dividing \eqref{eq:rank-one-sandwich} by \(N\), the remaining rank-one term is \(\sqrt{q_i/N}\to0\). Thus

\[
\frac1N\operatorname{tr}\sqrt{G_{N,i}}
\longrightarrow\sqrt{p_1(r)}.
\]

The sandwich \eqref{eq:trace-srm-bound} completes the SRM proof.
\end{proof}

\begin{remark}[An exact finite-size regime]\label{rem:exact-finite-regime}
If \(i\ge N-1\), then \(\min(|a-b|,i)=|a-b|\) for every matrix entry, and hence

\[
G_{N,i}=Q_N
\]

exactly. The approximation in Theorem~\ref{thm:long-known} is needed only when \(i<N-1\).
\end{remark}

\begin{figure*}[t]
  \centering
  \includegraphics[width=0.96\textwidth]{figures/figure2_analytic_limits.pdf}
  \caption{\textbf{Hierarchy of analytic localization limits.}
  (a) The fixed-length Toeplitz integral \(P_\infty(i,c^2)\) approaches the
  long-known-interval law \(p_1(c^2)\) as \(i\) grows. The \(i=2\) curve is
  equivalently given by the complete-elliptic-integral expression in
  Corollary~\ref{cor:two-site}. (b) An unknown interval instead yields
  \(p_1(c)^2\). The light-blue shaded band is the pointwise difference
  \(p_1(c^2)-p_1(c)^2\) between the two limiting laws; it is not a confidence
  band, a numerical error estimate, or a same-prior ignorance penalty.
  The one-translation/two-endpoint language is a structural interpretation;
  the second law follows from the two-dimensional Følner limit and the
  separable comparison symbol.}
  \label{fig:analytic-limits}
\end{figure*}

\section{Unknown interval length}\label{unknown-interval-length}

Let

\begin{align*}
\mathcal I_n&=\{[a,b]:1\le a\le b\le n\}, \\
M_n&=|\mathcal I_n|=\frac{n(n+1)}2,
\end{align*}

and equip \(\mathcal I_n\) with the uniform prior \(1/M_n\). This is uniform
over intervals, not over lengths. If \(L=|I|\), the induced marginal is

\begin{equation}\label{eq:induced-length-prior}
\begin{aligned}
\Pr\{L=\ell\}
&=\frac{n-\ell+1}{M_n}\\
&=\frac{2(n-\ell+1)}{n(n+1)},
\qquad 1\leq\ell\leq n.
\end{aligned}
\end{equation}

Conditionally on \(L=\ell\), the starting point is uniform over its
\(n-\ell+1\) admissible values. Consequently
\(\mathbb E L=(n+2)/3\), and \(L/n\) converges to the triangular density
\(2(1-x)\) on \(0<x<1\). Thus the prior favors shorter lengths, but its
typical interval still has order-\(n\) length. The candidate indexed by
\(I=[a,b]\) has \(|\psi\rangle\) on \(I\) and \(|0\rangle\) elsewhere. Its
exact Gram kernel is

\begin{equation}\label{eq:interval-gram}
G^{(n)}_{I,J}=c^{|I\triangle J|}.
\end{equation}

Theorem~\ref{thm:unknown-length} is therefore a Bayes law for the
uniform-over-intervals prior in \eqref{eq:induced-length-prior}, not a
prior-independent law for every unknown-length model. Alternative priors can,
for example, retain nonvanishing mass on bounded or logarithmically short
intervals, for which the reduction below need not apply.

For intersecting intervals,

\begin{equation}\label{eq:intersecting-distance}
|I\triangle J|=|a-a'|+|b-b'|.
\end{equation}

For disjoint intervals, assume without loss of generality that \(b<a'\) and define the separating gap by

\[
g(I,J):=a'-b-1\ge0.
\]

Then the exact identity is

\begin{equation}\label{eq:disjoint-distance}
|a-a'|+|b-b'|=|I\triangle J|+2g(I,J).
\end{equation}

Thus \eqref{eq:intersecting-distance} still holds for adjacent intervals, where \(g(I,J)=0\), but fails by twice the separating gap when \(g(I,J)>0\). The exact physical kernel is therefore not globally a two-dimensional Toeplitz kernel.

Classical multilevel Szeg{\H o} theorems describe spectral averages for genuine
multilevel Toeplitz compressions \cite{tilli1998,tyrtyshnikov1998}. Here the
triangular endpoint sets and, more importantly, the non-Toeplitz physical
kernel require the following local-diagonal Følner statement together with a
separate perturbative transfer back to \(G^{(n)}\).

\subsection{A Følner functional-calculus lemma}\label{a-folner-functional-calculus-lemma}

\begin{lemma}[Local diagonal Følner limit]\label{lem:folner-diagonal}
Let \(k\in\ell^1(\mathbb Z^d)\) satisfy \(k(-u)=\overline{k(u)}\), let \(A\) be convolution by \(k\) on \(\ell^2(\mathbb Z^d)\), and suppose its Fourier symbol satisfies

\[
0<\mu\le F(\theta)\le L<\infty
\]

almost everywhere. If \(\Lambda_n\subset\mathbb Z^d\) is any finite Følner sequence and \(A_n=P_nAP_n|_{\ell^2(\Lambda_n)}\), then for every continuous \(\phi:[\mu,L]\to\mathbb C\),

\begin{equation}\label{eq:folner-diagonal}
\frac1{|\Lambda_n|}\sum_{x\in\Lambda_n}
\left|
\langle\delta_x,\phi(A_n)\delta_x\rangle-a_\phi
\right|\longrightarrow0,
\end{equation}

where

\[
a_\phi=\frac1{(2\pi)^d}\int_{[0,2\pi]^d}\phi(F(\theta))\,d\theta.
\]
\end{lemma}

\begin{proof}
Since \(A\) has spectrum in \([\mu,L]\), the same is true of every compression
\(A_n\). Fix \(\varepsilon>0\) and put \(\eta=\varepsilon/8\). By the
Weierstrass theorem, applied to the real and imaginary parts if necessary,
choose a polynomial

\[
p(t)=\sum_{m=0}^{s}a_mt^m
\quad\text{such that}\quad
\sup_{t\in[\mu,L]}|\phi(t)-p(t)|<\eta.
\]

Let \(k_R=k\,\mathbf 1_{\{\|u\|_\infty\le R\}}\), let \(A_R\) be convolution
by \(k_R\), and set \(A_{R,n}=P_nA_RP_n|_{\ell^2(\Lambda_n)}\). Young's
inequality gives

\begin{equation}\label{eq:folner-young}
\|A-A_R\|_{\rm op},\ \|A_n-A_{R,n}\|_{\rm op}
\le\|k-k_R\|_{\ell^1}.
\end{equation}

Choose \(R\) so large that the right-hand side of
\eqref{eq:folner-young} is at most \(1\) and

\[
C_p\|k-k_R\|_{\ell^1}<\eta,\qquad
C_p:=\sum_{m=1}^{s}m|a_m|(L+1)^{m-1},
\]

where the second condition is void if \(C_p=0\). The telescoping identity for
\(X^m-Y^m\), together with
\(\|A\|,\|A_n\|\le L\) and
\(\|A_R\|,\|A_{R,n}\|\le L+1\), then yields

\begin{equation}\label{eq:folner-polynomial-truncation}
\|p(A)-p(A_R)\|_{\rm op},
\ \|p(A_n)-p(A_{R,n})\|_{\rm op}<\eta.
\end{equation}

Define

\[
Q_t=\{h\in\mathbb Z^d:\|h\|_\infty\le t\}
\]

and the \(sR\)-interior and boundary by

\begin{align*}
\Lambda_n^{\circ,sR}
&=\{x\in\Lambda_n:x+Q_{sR}\subseteq\Lambda_n\},\\
\partial_{sR}\Lambda_n
&=\Lambda_n\setminus\Lambda_n^{\circ,sR}.
\end{align*}

Every length-\(m\le s\) path generated by the range-\(R\) kernel and starting
at \(x\in\Lambda_n^{\circ,sR}\) stays in \(\Lambda_n\). Therefore

\begin{equation}\label{eq:folner-interior-equality}
\langle\delta_x,p(A_{R,n})\delta_x\rangle
=\langle\delta_0,p(A_R)\delta_0\rangle,
\qquad x\in\Lambda_n^{\circ,sR}.
\end{equation}

The radius here is \(sR\), not \(R\). Moreover,

\begin{equation}\label{eq:folner-boundary-inclusion}
\partial_{sR}\Lambda_n
\subseteq
\bigcup_{\|h\|_\infty\le sR}
\bigl(\Lambda_n\setminus(\Lambda_n-h)\bigr).
\end{equation}

The union is finite and each summand has \(o(|\Lambda_n|)\) elements by the
Følner property. Hence
\(|\partial_{sR}\Lambda_n|/|\Lambda_n|\to0\).
With
\(K_p=\sup_{|t|\le L+1}|p(t)|\), define

\begin{align*}
d_{n,R,p}(x)
&:=
\left|
\langle\delta_x,p(A_{R,n})\delta_x\rangle
-\langle\delta_0,p(A_R)\delta_0\rangle
\right|,\\
D_{n,R,p}
&:=\frac1{|\Lambda_n|}\sum_{x\in\Lambda_n}d_{n,R,p}(x).
\end{align*}

Equations
\eqref{eq:folner-interior-equality}--\eqref{eq:folner-boundary-inclusion} give

\begin{equation}\label{eq:folner-boundary-average}
D_{n,R,p}
\le2K_p\frac{|\partial_{sR}\Lambda_n|}{|\Lambda_n|}
\longrightarrow0.
\end{equation}

Finally, the spectral theorem and Fourier transformation identify

\[
\langle\delta_0,\phi(A)\delta_0\rangle
=\frac1{(2\pi)^d}\int_{[0,2\pi]^d}\phi(F(\theta))\,d\theta
=a_\phi.
\]

Put

\[
D_n^\phi:=\frac1{|\Lambda_n|}\sum_{x\in\Lambda_n}
\left|\langle\delta_x,\phi(A_n)\delta_x\rangle-a_\phi\right|.
\]

For all sufficiently large \(n\), functional calculus,
\eqref{eq:folner-polynomial-truncation}, and
\eqref{eq:folner-boundary-average} now give the single uniform error chain

\begin{align*}
D_n^\phi
&\le
\|\phi(A_n)-p(A_n)\|_{\rm op}\\
&\quad+\|p(A_n)-p(A_{R,n})\|_{\rm op}\\
&\quad+D_{n,R,p}\\
&\quad+\|p(A_R)-p(A)\|_{\rm op}\\
&\quad+\|p(A)-\phi(A)\|_{\rm op}\\
&<4\eta+\frac{\varepsilon}{2}=\varepsilon.
\end{align*}

This proves \eqref{eq:folner-diagonal}.
\end{proof}

\begin{lemma}[Gram upper bound for equiprobable pure states]\label{lem:gram-upper-bound}
Let \(G\) be an \(M\times M\) positive semidefinite Gram matrix with unit diagonal. Set

\[
d_j=(\sqrt G)_{jj},\qquad
q_j=\frac{d_j}{\operatorname{tr}\sqrt G},\qquad
u_j=\frac1M.
\]

Then

\begin{align}\label{eq:gram-upper-bound}
\left(\frac{\operatorname{tr}\sqrt G}{M}\right)^2
&\le P_{\rm opt}(G) \notag\\
&\le
\left(\frac{\operatorname{tr}\sqrt G}{M}\right)^2
+\sqrt{\lambda_{\max}(G)}\,\|q-u\|_1.
\end{align}
\end{lemma}

\begin{proof}
The lower bound is \eqref{eq:trace-srm-bound}. For \(G\succ0\), the upper
bound is Theorem 1 of Ref.~\cite{sentis2016}, rewritten in our normalization
so that no prior-weight or hypothesis-count factor is implicit. That theorem
assumes linearly independent states, equivalently a positive definite Gram
matrix. For singular \(G\), apply it to

\[
G_\varepsilon=(1-\varepsilon)G+\varepsilon I_M,
\qquad 0<\varepsilon<1.
\]

Each \(G_\varepsilon\) is positive definite and has unit diagonal. As \(\varepsilon\downarrow0\), the positive square root, its trace and diagonal, \(\lambda_{\max}\), and hence \(q\) all converge to their counterparts for \(G\). Lemma~\ref{lem:gram-stability} also gives \(P_{\rm opt}(G_\varepsilon)\to P_{\rm opt}(G)\). Passing to the limit proves the singular case.
\end{proof}

\subsection{The two-endpoint law}\label{the-two-endpoint-law}

\begin{theorem}[Unknown-length interval]\label{thm:unknown-length}
Under the uniform prior on \(\mathcal I_n\),

\begin{equation}\label{eq:unknown-length-limit}
P_{\rm opt}(G^{(n)}),P_{\rm SRM}(G^{(n)})
\longrightarrow p_1(c)^2
\end{equation}

for every fixed \(0\le c\le1\).
The convergence is pointwise in fixed \(c\). No finite-size rate or
uniformity as \(c\uparrow1\) is claimed.
\end{theorem}

\begin{proof}
For \(c=0\), the hypotheses are orthogonal, whereas for \(c=1\) they are
identical; both cases agree with the stated limit. Assume henceforth that
\(0<c<1\), and choose

\[
L_n=\lceil\gamma\log n\rceil,
\qquad \gamma>|\log c|^{-1},
\]

and retain only

\[
\mathcal B_n=\{I\in\mathcal I_n:|I|\ge L_n\}.
\]

The deleted prior mass is

\begin{equation}\label{eq:short-interval-mass}
\varepsilon_n
=1-\frac{|\mathcal B_n|}{M_n}
=O\!\left(\frac{\log n}{n}\right)\to0.
\end{equation}

On \(\mathcal B_n\), define the comparison kernel

\begin{equation}\label{eq:comparison-kernel}
H^{(n)}_{I,J}=c^{|a-a'|+|b-b'|}.
\end{equation}

It is the compression of the positive convolution kernel \(c^{|u|+|v|}\). The two kernels agree for intersecting intervals. For disjoint intervals,

\[
|G_{I,J}-H_{I,J}|
\le c^{|I|+|J|}\le c^{2L_n}.
\]

Consequently,

\begin{equation}\label{eq:comparison-operator}
\|G_{\mathcal B_n}-H_{\mathcal B_n}\|_{\rm op}
\le|\mathcal B_n|c^{2L_n}=o(1).
\end{equation}

Lemma~\ref{lem:gram-stability} transfers both success probabilities between these long-interval kernels.

We verify rather than assume the Følner property. Under the endpoint
identification \(I=[a,b]\leftrightarrow(a,b)\),

\[
\mathcal B_n
=\{(a,b)\in\mathbb Z^2:a\ge1,\ b\le n,\ b-a\ge L_n-1\},
\]

and

\begin{equation}\label{eq:long-domain-cardinality}
|\mathcal B_n|
=\frac{(n-L_n+1)(n-L_n+2)}2
\sim\frac{n^2}{2}.
\end{equation}

Fix \(h=(u,v)\in\mathbb Z^2\) and put \(H=\|h\|_\infty\). If the indicators
of \(\mathcal B_n\) at \(x=(a,b)\) and \(x+h\) differ, at least one of the
three affine inequalities defining \(\mathcal B_n\) changes truth value.
Such an \(x\) lies in a strip of respectively \(|u|\), \(|v|\), or
\(|v-u|\) lattice layers adjacent to
\(a=1\), \(b=n\), or \(b-a=L_n-1\). Because either \(x\) or \(x+h\) belongs
to \(\mathcal B_n\), each layer contains at most \(n+2H+1\) relevant points.
Put \(C_h=|u|+|v|+|v-u|\). Therefore

\begin{equation}\label{eq:long-domain-folner-bound}
\frac{|(\mathcal B_n-h)\triangle\mathcal B_n|}{|\mathcal B_n|}
\le
\frac{C_h(n+2H+1)}{|\mathcal B_n|}
\longrightarrow0,
\end{equation}

where \eqref{eq:long-domain-cardinality} and \(L_n=O(\log n)=o(n)\) were
used. Thus \((\mathcal B_n)\) is a Følner sequence in \(\mathbb Z^2\).
The comparison convolution symbol is

\[
F_c(\theta_1,\theta_2)=g_c(\theta_1)g_c(\theta_2),
\]

with a uniform positive lower bound. Lemma~\ref{lem:folner-diagonal} with \(\phi(t)=\sqrt t\) yields

\begin{equation}\label{eq:folner-sqrt-diagonal}
\frac1{|\mathcal B_n|}\sum_{I\in\mathcal B_n}
\left|(\sqrt{H_{\mathcal B_n}})_{II}-\alpha_c\right|
\to0,
\end{equation}

where separability gives

\begin{align}\label{eq:endpoint-factorization}
\alpha_c
&=\frac1{(2\pi)^2}\int\!\!\int\sqrt{F_c} \notag\\
&=\left[\frac1{2\pi}\int_0^{2\pi}
\sqrt{g_c(\theta)}\,d\theta\right]^2
=p_1(c).
\end{align}

The lower bound \eqref{eq:trace-srm-bound} therefore converges to \(\alpha_c^2\). For the matching upper bound, put

\[
q_I=\frac{(\sqrt H)_{II}}{\operatorname{tr}\sqrt H},
\qquad u_I=\frac1{|\mathcal B_n|}.
\]

Equation~\eqref{eq:folner-sqrt-diagonal} implies \(\|q-u\|_1\to0\), and

\[
\lambda_{\max}(H_{\mathcal B_n})
\le\|F_c\|_\infty
=\left(\frac{1+c}{1-c}\right)^2.
\]

Lemma~\ref{lem:gram-upper-bound} gives

\[
P_{\rm opt}(H_{\mathcal B_n})
\le
\left(\frac{\operatorname{tr}\sqrt H}{|\mathcal B_n|}\right)^2
+\sqrt{\lambda_{\max}(H)}\,\|q-u\|_1,
\]

so the optimum tends to \(\alpha_c^2=p_1(c)^2\). Equation~\eqref{eq:folner-sqrt-diagonal} also implies convergence of the average squared diagonal and hence of the SRM.

We now pass from the long-interval subfamily back to all intervals. Deleting
prior mass \(\varepsilon_n\) changes the optimum by at most
\(\varepsilon_n\), which already proves the full-family optimal limit. The SRM
needs an additional trace comparison. Set

\[
t_n=\frac{\operatorname{tr}\sqrt{G^{(n)}}}{M_n}.
\]

Since \(G_{\mathcal B_n}\) is a principal submatrix, eigenvalue interlacing gives

\begin{equation}\label{eq:full-trace-interlacing}
t_n\ge
\frac{|\mathcal B_n|}{M_n}
\frac{\operatorname{tr}\sqrt{G_{\mathcal B_n}}}{|\mathcal B_n|}.
\end{equation}

The factor \(|\mathcal B_n|/M_n=1-\varepsilon_n\) is essential. Equations~\eqref{eq:comparison-operator}--\eqref{eq:endpoint-factorization} show that the right-hand side tends to \(\alpha_c\), whereas \(t_n^2\le P_{\rm opt}(G^{(n)})\) gives the reverse limsup. Thus \(t_n\to\alpha_c\), and

\[
t_n^2\le P_{\rm SRM}(G^{(n)})\le P_{\rm opt}(G^{(n)})
\]

proves the original, untruncated SRM limit.
\end{proof}

\section{Joint detection and exact localization}\label{joint-detection-and-exact-localization}

Add the no-anomaly state

\[
|\Omega\rangle=|0\rangle^{\otimes n}
\]

with fixed prior \(\pi_0\in(0,1)\). The task is successful only if the measurement reports \(H_0\) when no anomaly is present, or reports the exact anomalous interval otherwise.

\begin{theorem}[Fixed known length with no-anomaly prior]\label{thm:h0-known-length}
For fixed \(i\),

\begin{equation}\label{eq:h0-fixed-limit}
P_{\rm joint}^*(N,i,r,\pi_0)
\longrightarrow
\pi_0+(1-\pi_0)P_\infty(i,r).
\end{equation}

If \(i=i_n\to\infty\) and \(N_n\to\infty\), then

\begin{equation}\label{eq:h0-growing-limit}
P_{\rm joint}^*
\longrightarrow
\pi_0+(1-\pi_0)p_1(c^2).
\end{equation}
\end{theorem}

\begin{proof}
If \(r=1\), every anomalous candidate equals \(|\Omega\rangle\). The \(N+1\)
labels are therefore represented by one physical state, with label priors
\(\pi_0\) and \((1-\pi_0)/N\). Hence, exactly at finite \(N\),

\begin{equation}\label{eq:h0-fixed-r-one}
P_{\rm joint}^*(N,i,1,\pi_0)
=\max\!\left\{\pi_0,\frac{1-\pi_0}{N}\right\}
\longrightarrow\pi_0.
\end{equation}

This agrees with \eqref{eq:h0-fixed-limit} because
\(P_\infty(i,1)=0\). Assume henceforth that \(0\le r<1\).
Put \(q=r^i\) and use

\[
|\Phi_a^{(i)}\rangle
=\sqrt q|\Omega\rangle+\sqrt{1-q}|\xi_a^{(i)}\rangle.
\]

If \(P_G^*\) and \(P_R^*\) are the equal-prior localization optima for the original and residual ensembles, respectively, then the explicit measurement that assigns \(|\Omega\rangle\) to \(H_0\) and uses an optimal residual measurement on \(|\Omega\rangle^\perp\) gives

\[
\pi_0+(1-\pi_0)(1-q)P_R^*
\le P_{\rm joint}^*.
\]

Conversely, the no-anomaly contribution is at most \(\pi_0\), and completing the anomalous sub-POVM can only improve its conditional localization score, so

\[
P_{\rm joint}^*\le\pi_0+(1-\pi_0)P_G^*.
\]

For fixed \(i\), Corollary~\ref{cor:residual-toeplitz} gives
\((1-r^i)P_R^*\to P_\infty(i,r)\), and the two bounds prove
\eqref{eq:h0-fixed-limit}. In the growing-length regime,
\(|P_G^*-P_R^*|\le r^{i/2}\) and \(q\to0\); Theorem~\ref{thm:long-known}
then proves \eqref{eq:h0-growing-limit}. At \(r=0\), all \(N+1\) states are
mutually orthogonal and \(P_{\rm joint}^*=1\), also in agreement with both
displayed limits.
\end{proof}

\begin{theorem}[Unknown length with no-anomaly prior]\label{thm:h0-unknown-length}
Under the uniform prior over all nonempty intervals,

\begin{equation}\label{eq:h0-unknown-limit}
P_{\rm joint}^*(n,c,\pi_0)
\longrightarrow
\pi_0+(1-\pi_0)p_1(c)^2.
\end{equation}
\end{theorem}

\begin{proof}
Assume \(0<c<1\) and use the long-interval family \(\mathcal B_n\) from Theorem~\ref{thm:unknown-length}. For \(I\in\mathcal B_n\), put

\[
s_I=\langle\Omega|\Phi_I\rangle=c^{|I|}\le c^{L_n}=:\eta_n\to0
\]

and

\[
|\xi_I\rangle
=\frac{|\Phi_I\rangle-s_I|\Omega\rangle}{\sqrt{1-s_I^2}}.
\]

To make the conditioning explicit, define the full-family anomalous
localization optimum

\[
P_{G,n}^{\rm all}
:=\max_{\{M_I\}_{I\in\mathcal I_n}}
\frac1{M_n}\sum_{I\in\mathcal I_n}
\langle\Phi_I|M_I|\Phi_I\rangle
=P_{\rm opt}(G^{(n)}),
\]

and define the two equal-prior optima on the retained family by

\begin{align*}
P_{G,n}^{\mathcal B}
&:=\max_{\{M_I\}_{I\in\mathcal B_n}}
\frac1{|\mathcal B_n|}\sum_{I\in\mathcal B_n}
\langle\Phi_I|M_I|\Phi_I\rangle,\\
P_{R,n}^{\mathcal B}
&:=\max_{\{M_I\}_{I\in\mathcal B_n}}
\frac1{|\mathcal B_n|}\sum_{I\in\mathcal B_n}
\langle\xi_I|M_I|\xi_I\rangle.
\end{align*}

In each maximization the operators form a POVM on the support of the
corresponding conditional ensemble, completed arbitrarily on its orthogonal
complement. These are conditional localization probabilities: the factors
\(1-\pi_0\) and \(1-\varepsilon_n\) are not included in their definitions.
The pure-state trace distance gives the uniform comparison

\[
|P_{R,n}^{\mathcal B}-P_{G,n}^{\mathcal B}|\le\eta_n.
\]

Use the optimal residual POVM for the long intervals in \(|\Omega\rangle^\perp\), assign zero operators to deleted short intervals, and assign the remaining operator to \(H_0\). With \(\varepsilon_n=1-|\mathcal B_n|/M_n\), this gives

\[
P_{\rm joint}^*
\ge\pi_0+(1-\pi_0)(1-\varepsilon_n)(1-\eta_n^2)
(P_{G,n}^{\mathcal B}-\eta_n).
\]

The universal upper bound is

\[
P_{\rm joint}^*
\le\pi_0+(1-\pi_0)P_{G,n}^{\rm all}.
\]

Theorem~\ref{thm:unknown-length} and \(\varepsilon_n,\eta_n\to0\) make the bounds coincide. If \(c=0\), all hypotheses are mutually orthogonal and the success probability is one. If \(c=1\), all hypotheses are identical and the optimum is \(\max\{\pi_0,(1-\pi_0)/M_n\}\to\pi_0\).
\end{proof}

\section{Numerical certification}\label{numerical-certification}

The numerical calculations verify the finite-size Gram matrices and show how
the asymptotic regimes are approached. None of the proofs relies on numerical
evidence.

For small unknown-length instances, we solve the minimum-error problem in
canonical Gram coordinates with both primal and dual semidefinite programs.
The certification data cover \(n=3,\ldots,7\) and

\[
c\in\{0.3,0.6,0.8,0.9,0.95,0.99\}.
\]

Raw solver objectives are retained as diagnostics but are not used directly as
bounds. Let \(\rho_j\) denote the prior-weighted canonical states on their
\(d_{\rm supp}\)-dimensional support. For a numerical dual matrix \(Y\), set

\[
\delta_D=\max_j\max\{0,-\lambda_{\min}(Y-\rho_j)\}
\]

and add the recorded scale-aware roundoff margin \(\zeta_D>0\). Then

\begin{equation}\label{eq:safe-dual-bound}
Y_{\rm safe}=Y+(\delta_D+\zeta_D)I,\qquad
U_{\rm safe}=\operatorname{tr}Y_{\rm safe}
\end{equation}

is rechecked against every dual PSD constraint and supplies a feasible upper
bound.

For the primal matrices, project each raw \(M_j\) onto the PSD cone, write the
result as \(M_j^+=B_jB_j^*\), and put \(S=\sum_jM_j^+\). The normalized
factors \(S^{-1/2}B_j\) give PSD operators summing to the identity in exact
arithmetic. In floating point, the implementation additionally contracts their
sum \(T\) by

\[
\gamma=\bigl(\max\{1,\lambda_{\max}(T)\}+\zeta_P\bigr)^{-1}
\]

and assigns the PSD remainder \(I-\gamma T\) to one outcome. The resulting
POVM \(\{\widehat M_j\}\) is rechecked and gives

\begin{equation}\label{eq:safe-primal-bound}
L_{\rm safe}=\sum_j\operatorname{tr}(\rho_j\widehat M_j),\qquad
L_{\rm safe}\le P_{\rm opt}\le U_{\rm safe}.
\end{equation}

Across the 30 unknown-length cases, the largest raw absolute primal--dual gap
is \(6.782\times10^{-9}\), while the largest postprocessed feasible interval
width \(U_{\rm safe}-L_{\rm safe}\) is \(8.854\times10^{-9}\). The largest
safe completeness residual is \(1.38\times10^{-15}\), and every rechecked
safe primal and dual PSD violation is zero. The smallest safe primal
eigenvalue and dual slack eigenvalue are respectively
\(3.4103\times10^{-13}\) and \(3.4102\times10^{-13}\). Full dense
physical-Gram SRM computations for the single-interval
model extend to \(n=50\), with candidate count \(M_n=1275\), over

\[
c\in\{0.5,0.8,0.9,0.95,0.99\}.
\]

High-overlap cases converge slowly. The correlation length

\[
\xi(c)=|\log c|^{-1}
\]

is approximately \(99.5\) at \(c=0.99\), so even \(n=50\) has \(n/\xi\approx0.503\). The corresponding finite-size deviations should not be interpreted as asymptotic counterexamples.

\begin{figure*}[t]
  \centering
  \includegraphics[width=0.92\textwidth]{figures/figure3_finite_size.pdf}
  \caption{\textbf{Finite-size convergence and certified SRM-gap intervals.}
  (a) SRM success for \(i=\lceil\sqrt N\rceil\) and
  \(i=\lfloor N/2\rfloor\); dotted horizontal lines are \(p_1(c^2)\).
  (b) Full dense physical-Gram SRM values through \(n=50\); dotted lines are
  \(p_1(c)^2\), with high-overlap curves retained as stress tests.
  (c) Postprocessed feasible brackets:
  \(L_{\rm safe}-P_{\rm SRM}\le P_{\rm opt}-P_{\rm SRM}
  \le U_{\rm safe}-P_{\rm SRM}\) for \(n\leq7\). Plotting symbols mark only
  the midpoint of each displayed bracket, not an estimate of
  \(P_{\rm opt}-P_{\rm SRM}\); error bars are often smaller than the symbols.}
  \label{fig:finite-size}
\end{figure*}

The accompanying reproducibility package records the commands, package
versions, file hashes, solver residuals, and data files. The environment is
pinned to Python 3.12 and fixed versions of NumPy, SciPy, Matplotlib, CVXPY,
Clarabel, and SCS; the reported clean-environment rerun used Python 3.12.10.

\section{Discussion and limitations}\label{discussion-and-limitations}

Known and unknown lengths produce different analytic laws. A long known interval is indexed by one translation coordinate. Moving it by one site changes two local factors, hence the elementary kernel parameter is \(c^2\) and the limit is \(p_1(c^2)\). An unknown interval is indexed by two endpoints. In the long intersecting regime, the symmetric-difference distance separates into the \(\ell^1\) distance of these endpoints; each endpoint contributes a kernel with parameter \(c\). The normalized square-root trace is therefore \(p_1(c)\), and the success probability is \(p_1(c)^2\). This factorized law depends on the uniform-over-intervals prior: it should not be read as a universal penalty for every prior on an unknown duration, nor as a same-prior value-of-information comparison with the known-length model.

For fixed known length, Theorem~\ref{thm:fixed-length} proves the rigorous rate \(P_{\rm opt}-P_{\rm SRM}=O_{i,r}(N^{-1/2})\). Exploratory finite-size calculations in some fixed-length parameter ranges are consistent with an \(N^{-1}\) decay, but the available certified sizes do not establish a uniform exponent. Determining whether the wrap-around boundary perturbation can be exploited to prove \(O_{i,r}(N^{-1})\), or an asymptotic expansion with a nonzero \(N^{-1}\) coefficient, remains an open problem. By contrast, Theorems~\ref{thm:long-known} and \ref{thm:unknown-length} are qualitative fixed-overlap limits: their proofs give neither a finite-size rate nor bounds uniform in a local regime \(c=c_n\uparrow1\).

The theorems benchmark arbitrary collective POVMs; they do not provide an
efficient physical implementation of the optimum measurement. Nor do the
known collective advantages for permanent changes establish an advantage for
this returning-interval prior: a matched local or sequential baseline remains
an open operational comparison. The results also rely
on a known pure anomalous state, a fixed overlap \(c\), and uniform priors. In
particular, the analysis does not cover a local asymptotic regime
\(c=c_n\to1\), unknown or mixed anomalous states, minimax priors, restricted
local measurements, or a growing number of disjoint intervals. Each extension
changes either the hypothesis geometry or the admissible measurement class
and requires a separate analysis.

\section{Conclusion}\label{conclusion}

A returning quantum change retains enough translation structure for an exact
asymptotic analysis. With fixed known length, the localization probability is
a Toeplitz square-root integral; a growing known length gives the
single-change law \(p_1(c^2)\), whereas an unknown length gives the
two-endpoint law \(p_1(c)^2\). The SRM reaches the global minimum-error limit
in all three regimes. A no-change hypothesis with fixed prior adds its Bayes
contribution without altering the conditional localization law. Thus the
contiguity constraint turns a general quantum multiple-hypothesis problem
into tractable one- and two-dimensional spectral limits.

\section*{Data and code availability}

This work is theoretical and generated no experimental data. All code and
numerical data required to reproduce the results are available at
\url{https://github.com/chenxu8817-cyber/quantum-change-intervals}. The
immutable version corresponding to this manuscript is GitHub release
\href{https://github.com/chenxu8817-cyber/quantum-change-intervals/releases/tag/v1.0.0-paper1}{v1.0.0-paper1}.
An archival DOI will be added if that release is deposited in a permanent
repository.

\section*{Author contributions}

Xu Chen: Conceptualization; development of the physical model; formal
analysis; mathematical derivation and proof of the theorems.

Xue Ma: Numerical computation; software implementation; generation of
numerical data; data visualization and figure preparation.

Both authors contributed to the interpretation of the results, writing, and
revision of the manuscript.

\section*{AI-assisted tools}

The authors used OpenAI Codex as an assistive tool for literature discovery
and bibliographic checking, LaTeX formatting, code and numerical
cross-checks. The authors independently verified every cited source, theorem,
derivation, numerical output, figure, and final statement and take full
responsibility for the work.

\section*{Funding}

This work was supported by Hebei University of Science and Technology under
Starting Grant No.~81/1181298.

\section*{Competing interests}

The authors declare no competing interests.


